\documentclass[10pt]{article}
\usepackage[margin=20mm]{geometry}
\usepackage{amsmath,amssymb,mathtools,fontspec,unicode-math,xurl,hyperref}
\IfFontExistsTF{DejaVu Serif}{\setmainfont{DejaVu Serif}}{\setmainfont{Latin Modern Roman}}
\setmathfont{latinmodern-math.otf}
\setlength{\parindent}{0pt}\setlength{\parskip}{6pt}\sloppy
\title{Scientific initial review of M25-0006
On approximation to a real number by algebraic numbers of bounded degree}
\author{}
\date{}
\begin{document}
\maketitle
=3em\par
\section*{Identity, scope, and disposition}
The reviewed object is Anthony Poëls’s arXiv:2405.08341v1, dated 14 May 2024, comprising 21 PDF pages. The frozen journal identity is Annals of Mathematics 201 (2025), 307–330, DOI `10.4007/annals.2025.201.1.6`. These identities correspond, but equivalence to the journal Version of Record has not been established. All seven manifest file bindings passed the actual byte-count and SHA-256 checks. The reviewed PDF SHA-256 is\par
00001TOKEN.\par
All 21 pages were read. Definitions, the quantifiers in the approximation exponents, the polynomial construction, generalized resultants, the finite-dimensional minimization and final constant were checked. Pages 1, 13, 14, 16, 19 and 20 were also inspected as rendered source images. The 1,367-line bound author TeX was parsed statically; it was not executed. There are no active external input commands in that source.\par
Disposition: one Minor finding, restricted to the v1 abstract; no demonstrated Major defect in the reviewed core argument. Several compressed steps and local formula slips are resolved explicitly below. This is a scientific initial review with the stated coverage, not a whole-paper correctness certificate or verification of every cited theorem. In particular, the positive conclusion is conditional on the classical geometry-of-numbers and transference results identified below.\par
\section*{F01: the abstract defines the exponent using the wrong extremum}
Severity and location. Minor. PDF page 1, abstract, third sentence; `source.tex`, line 326. The abstract uses infimum where the body uses supremum, in the definition of \(\omega_n^*(\xi)\).\par
Equation reference: eq:down. Complete check. Let \(E\) be the set of exponents \(w\) for which infinitely many algebraic numbers of degree at most \(n\) satisfy \[
 0<|\xi-\alpha|\le H(\alpha)^{-w-1}.\] For \(H(\alpha)\ge1\), if \(w\in E\) and \(v<w\), the same solutions satisfy the inequality for \(v\). Thus \(E\) is downward closed. It is nonempty: for irrational \(\xi\), rational convergents have \(|\xi-p/q|\le q^{-2}\) and \(H(p/q)=\max(|p|,q)\ll_\xi q\), which gives every fixed \(0<w<1\) after discarding finitely many convergents. Consequently its infimum is \(-\infty\) if all real exponents are allowed, or \(0\) if the domain is restricted to positive exponents. Neither is the positive exponent claimed in the abstract. The supremum has the required meaning.\par
Counterevidence and impact. Equation (1.2) on the same page and Section 2 on page 4 correctly specify the supremum. The construction and Theorem 6.1 use arbitrarily large heights and then let \(\varepsilon\) decrease to zero in this correct sense. The defect is local to the abstract and does not invalidate Theorem 1.1. Replacing the single extremum word repairs it.\par
Version attribution. The current publisher abstract already uses the correct extremum. The publisher page records a revision dated 19 July 2024, after v1. This verifies correction in the publisher’s abstract only; it does not verify the entire journal PDF or show exactly when the correction was made. F01 is not alleged to remain in the journal article.\par
\section*{Polynomial construction and dependency checks}
\subsection*{Quantifiers, minima, and irreducible factors}
The distinction between a solution for every sufficiently large height and one for arbitrarily large heights is maintained. The reduction \(\omega_n(\xi)=\infty\Rightarrow\omega_n^*(\xi)=\infty\) uses the stated classical transference inequality; the finite case permits the uniform lower bound (4.3) after finitely many small polynomials are excluded. Transcendence ensures that a nonzero integer polynomial never vanishes at \(\xi\) and that a bounded-height collection is finite.\par
For the convex body on page 4, its volume is comparable to \(e^{-q}\) with constants depending on \(n,\xi\): the coefficient of the constant term in evaluation is 1, and fixed smaller boxes give the lower bound. Minkowski’s second theorem therefore gives \(\sum L_i(q)=q+O(1)\). The unboundedness of \(L_1\) follows by finiteness of integer polynomials of bounded height. In Lemma 3.1 the strict height inequality can be obtained by first choosing an exponent strictly between the working exponent and the uniform exponent, then using a fixed factor smaller than the target height. The exponent margin absorbs this factor.\par
For Lemma 3.3, changing signs of the initial \(Q_i\) so that their values have the sign of \(Q_1(\xi)\) makes the floor subtraction immediately satisfy the displayed factor-2 norm estimate. Such sign changes preserve successive-minimum realization and independence. Subtracting multiples of \(Q_1\) preserves independence; sorting the resulting heights preserves the upper bounds by the ordered successive minima. The integer determinant expanded in its evaluation row supplies the lower product bound, while Minkowski supplies the upper bound. Without the harmless sign convention, the printed elementary estimate needs a larger absolute constant; this does not change any later exponent.\par
The infinitude of the irreducible family in Section 4 uses a standard factor argument. Factor sufficiently good integer polynomials into their content and at most \(n\) irreducible factors. Product-height comparability and an exponent margin imply that at least one factor has the required small value. A finite collection of possible factors, with multiplicities bounded by \(n\), cannot account for values tending to zero with unbounded heights. Constants are absorbed by the fixed positive exponent margin. The invoked transference and product-height theorems are retained as external inputs, rather than presented as newly proved here.\par
\subsection*{The two branches of Lemma 4.1}
The maximal \(q\) with \(L_1(q)=\log(e^{-n}\lVert P\rVert)\) exists for the large heights in use. Its terminal trajectory has slope 1. Gelfond’s inequality and irreducibility give coprimality of \(P\) and \(Q\), including the possibility that their degrees are smaller than \(n\). Adding \(3P\) only to the low-height basis vectors makes their heights exceed \(\lVert P\rVert\) while preserving spanning together with \(P\). A basis-exchange argument then retains \(Q,P\) and removes one of the other vectors.\par
In the \(P_0\) branch, \(q<q_P\) implies \(e^{-n}|P(\xi)|\le|Q(\xi)|\); in the other branch, \(q\ge q_P\) implies \(|Q(\xi)|<|P(\xi)|\). Together with the determinant lower bound these yield respectively the exponent ranges around \(\widehat\omega_n\) and \(\omega_n\). Their union supplies the weaker uniform estimate needed in Theorem 6.1. The constants depending on \(n,\xi\) are fixed before the height tends to infinity; thresholds may depend on \(\varepsilon\).\par
\section*{Generalized resultants, including the endpoint height branch}
\subsection*{Dimension and root estimates}
The cited dependency [13], On uniform polynomial approximation, arXiv:2405.07219v1, Section 6, was consulted directly. Its Lemma 6.3 and Proposition 6.2 have the required hypotheses. The essential argument can also be reconstructed: writing \(V_i=V_i(\mathcal A)\), one has \(V_{i+1}=V_i+XV_i\). The quotient map from \(V_i\) to \(V_{i+1}/XV_i\) is onto and its kernel contains \(XV_{i-1}\). Hence the dimension increments decrease. Strict increase follows from a shifted nonzero polynomial of degree \(i+1\). If \(\gcd(\mathcal A)=1\), a Bezout identity puts 1 in some \(V_\ell\), and sufficiently high spaces are full. Concavity then gives the dimension bound used in Lemma 5.2. For the actual sets containing a coprime pair, full dimension at \(2n-1\) also follows directly by the Sylvester-map injection.\par
The degree-padding factors \((X-\lambda_i)^{n-\deg P_i}\) can be chosen with \(\lambda_i\in\{0,\ldots,n\}\) avoiding the other polynomial’s roots and each other when necessary. Thus they preserve coprimality and have controlled heights and values. The basis extension has the asserted dimension. Taylor translation has determinant 1, so the generalized resultant remains a nonzero integer after translation.\par
For any nonconstant polynomial \(S\) and nearest complex root \(\alpha\) to \(\xi\), \[\left|\frac{S'(\xi)}{S(\xi)}\right|
 \le \sum_{\beta:S(\beta)=0}\frac1{|\xi-\beta|}
 \le \frac{\deg S}{|\xi-\alpha|}.\] This proves the root-distance estimate including repeated roots. Every root in the constructed basis has degree at most \(n\): added roots are integers or zero, and the others come from degree-at-most-\(n\) polynomials. The statement uses algebraic numbers without a real-root restriction; no unproved assertion that the nearest root is real is needed.\par
\subsection*{Why the index \(m=1\) may be removed in Corollary 5.4}
The passage from Proposition 5.1 to (5.6), page 14, is compressed: the proposition permits \(m=1\), whereas the corollary writes \(m\ge2\). It requires justification, especially at \(k=n+1\). Here is a direct completion, with constants independent of \(C,y\).\par
Set \(r=n-k+2\). If \(k\le n\) and the root in Proposition 5.1 comes from the first-height group, its distance is bounded by \[\delta^2 H_1^{r-2}H_2^r H_3\cdots H_k
 \le \delta^2 H_2^{2r-2}H_3\cdots H_k.\] Its height is at most a fixed multiple of \(H_1\le H_2\), so this is precisely the required estimate with height index 2. Other indices work directly.\par
Equation reference: eq:refine. For \(k=n+1\), use the first two rows of the determinant expansion before discarding height information. Give its \(k\) columns height budgets \(B_1=H_1,B_2=H_2,\ldots,B_k=H_k\), allowing fixed factors from degree padding. Even if an actual column has smaller height, these budgets remain valid. If the derivative is in column 1, the value row must use another column with budget at least \(H_2\). If it is in another column, the remaining smallest budget is at least \(H_1\). Consequently, with \(M_i=\max(B_i,H_2)\), \[
 1\ll \delta\,\frac{\prod_i B_i}{H_1}
          \sum_i\frac{|S_i'(\xi)|}{M_i}.\] Choose a largest summand, and apply the nearest-root estimate. It gives \[|\xi-\alpha|\ll
 \delta^2\,\frac{H_2\cdots H_k}{M_i},
 \qquad H(\alpha)\ll M_i.\] Here \(M_i\) is one of \(H_2,\ldots,H_k\), exactly as (5.6) requires. This resolves the endpoint, rather than assuming \(H_1\) and \(H_2\) are comparable.\par
The sentence after (5.4), page 13, writes \(|S_\ell'(\xi)|/\lVert S_\ell\rVert^{-1}\) (source line 876). The largest summand actually needed is \(|S_\ell'(\xi)|/\lVert S_\ell\rVert\), as displayed in (5.4). This local selection typo and the omitted branch explanation have a complete direct resolution above; they are not evidence that the approximation bound or main theorem is false.\par
With \(A_k\ge0\), conversion from a height budget to \(H(\alpha)\) has the correct inequality direction. In Theorem 6.1, \(A_k(x,a)\ge1\) and \(\varepsilon<1/2\) ensure this sign condition. There are finitely many \(k\), so maximizing over \(k\) does not lose uniformity. The independent bound \(|\xi-\alpha|\ll H_2^{-1}\) forces \(H(\alpha)\to\infty\), excluding recycling finitely many roots. The exponent margin also absorbs the fixed constant in the definition of \(\omega_n^*\).\par
\section*{Minimization and boundary analysis}
For \(x=n\), the feasible vector is \((1,\ldots,1)\) and \(F=n\). For \(x\ge(3n-1)/2\), the \(k=2\) term is at least \(n\). The trial vector \((1,\ldots,1,1+\varepsilon)\), \(0<\varepsilon<1/2\), has every ratio below \(n\). The page-16 display (source line 1028) should label its left side \(A_{n+1}/a_{n+1}\): the numerator itself is \(n+\varepsilon\), and the required ratio is \((n+\varepsilon)/(1+\varepsilon)\). This local label slip does not prevent the stated strict bound for \(F\). Compactness yields a minimizer with strict interior \(x\) bounds.\par
Reducing one coordinate immediately after a strict jump lowers all other ratios. At a minimizer this forces the jump ratio to be active. The same fact rules out coordinates at least 2 and equal consecutive coordinates above the initial plateau. This proves the active-set structure in Lemma 7.2 without assuming differentiability of a maximum.\par
The recurrence for the remaining coordinates is \[a_{k+1}=\frac{F}{F+1}a_k+\frac2{F+1}.\] Its explicit solution gives the sum constraint (7.6). At \(\theta=1/2\), that constraint makes \(F=(2n-\ell+1)/2\) rational. Writing \(F/(F+1)=a/b\) in lowest terms, \(2a^m=b^m\) with \(m=n+1-\ell\) forces \(m=1,a=1,b=2\). Thus \(F=1,\ell=n\) and \(n=1\), contrary to \(n\ge2\). Division by \(1-2\theta\) in the perturbation argument is therefore legitimate. The strict bounds persist for a small decrease of \(F\), contradicting minimality unless \(F=2(x-n)+\ell-1\).\par
Finally, the two minimizing identities give \[(3-2\theta)F_n=2n+1-\ell,\qquad
 F_n\log\theta\ge-(n+1-\ell).\] The second inequality follows from \(\log(1+u)\le u\), \(u>0\), with the sign checked after multiplication by \(-(n+1-\ell)\). Adding them yields \((3-2\theta+\log\theta)F_n\ge n\). The scalar function has derivative \(-2+1/\theta\) and negative second derivative, hence global maximum \(2-\log2\) at \(\theta=1/2\). This proves the final lower bound. The exceptional value excluded for an exact minimizer is used only as a global upper bound on this scalar function, which is logically consistent.\par
\section*{Conclusion and remaining decisions}
Item | Judgment | Repair, adoption, or remaining action\par
F01, v1 abstract | Minor; incorrect extremum | Replace infimum by supremum. The body already uses the correct definition; the publisher abstract also has the correction. Author adoption was not requested.\par
Core approximation bound | No demonstrated Major defect found | Endpoint height selection and optimization steps have explicit completions above. Classical transference and geometry-of-numbers inputs remain external dependencies.\par
Version scope | Formal review of arXiv v1 | Only the publisher abstract was directly compared. The full journal PDF and full version equivalence remain unverified.\par
Decision | Initial review complete | Scientific/editorial disposition remains pending. No whole-paper correctness certificate or author contact.\par
\section*{Supporting checks}
The assistant’s own standard-library Python executed successfully: 60 exact rational Taylor-determinant identities; 119 exact generalized-resultant rank checks on explicit coprime families; 55 exact recurrence/sum identities; 165 exact height-budget inequalities; and scalar optimization checks for all \(2\le n\le64\). It gives \(F_2=\sqrt3\) numerically and \(F_{64}\simeq49.2350063\), each consistent with the analytic bound. Finite checks are supporting evidence, not proofs for all degrees. The script SHA-256 is\par
00002TOKEN.\par
The live arXiv record still lists only v1. The publisher records receipt on 16 May 2024, revision on 19 July, acceptance on 6 September, and online publication on 9 January 2025. A bounded title/correction/erratum search did not identify a separate article-specific correction. Absence from that search is not proof of absence. Only the abstract correction is directly compared; the journal final remains unacquired.\par
This exact English TeX is compiled by the previously authorized Tectonic binary under untrusted mode, without shell escape or author-code execution. Completion is recorded only after the resulting PDF has been rendered and every page inspected.\par
3 A. Poëls, On approximation to a real number by algebraic numbers of bounded degree, arXiv:2405.08341v1 (14 May 2024). Bound input PDF and author TeX. https://arxiv.org/abs/2405.08341. Annals of Mathematics, article metadata and abstract, DOI 10.4007/annals.2025.201.1.6, consulted during this review. https://annals.math.princeton.edu/2025/201-1/p06. A. Poëls, On uniform polynomial approximation, arXiv:2405.07219v1, Section 6, especially Lemma 6.3 and Proposition 6.2. https://arxiv.org/html/2405.07219v1.\par
\end{document}
